\documentclass[../../main.tex]{subfiles}

\begin{document}

\chapter{Struktur Kontrol dan Percabangan}
\label{chap:kontrol}

% =========================================================
% OUTCOME PEMBELAJARAN KHUSUS BAB
% =========================================================

\begin{subcpmkbox}
\SubCPMK{3.3} --- Mengimplementasikan struktur kontrol (\asm{CMP},
\asm{TEST}, lompatan bersyarat \asm{Jcc}, \asm{LOOP}, dan prosedur
dengan \asm{CALL}/\asm{RET}) dalam program assembly.

\vspace{0.4em}
\textbf{CPMK yang didukung:} \CPMK{03} \quad\textbullet\quad
\textbf{CPL yang didukung:} \CPL{03}
\end{subcpmkbox}

\begin{tujuanbox}
Setelah mempelajari bab ini, mahasiswa mampu:
\begin{enumerate}
    \item menjelaskan peran \eng{flag} sebagai dasar pengambilan
          keputusan;
    \item membedakan cara kerja \asm{CMP} dan \asm{TEST};
    \item memilih \asm{Jcc} yang tepat untuk bilangan bertanda dan
          tak bertanda;
    \item menjelaskan pengaruh jarak lompatan terhadap ukuran kode;
    \item menyusun perulangan dengan \asm{LOOP}, \asm{LOOPE}, dan
          \asm{LOOPNE} serta membandingkannya dengan pola
          \asm{DEC} dan \asm{JNZ};
    \item menerjemahkan pola \code{if}, \code{if--else}, \code{while},
          \code{do--while}, \code{for}, dan \code{switch} dari bahasa
          tingkat tinggi ke assembly;
    \item menuliskan program pencarian nilai maksimum dan pengurutan
          sederhana menggunakan struktur kontrol;
    \item menghindari kekeliruan umum pada label dan lompatan.
\end{enumerate}
\end{tujuanbox}

\begin{materibox}
\begin{enumerate}
    \item Konsep \eng{flag} sebagai dasar pengambilan keputusan;
          kaitan dengan \flag{ZF}, \flag{CF}, \flag{SF}, \flag{OF},
          dan \flag{PF}.
    \item \asm{CMP} dan \asm{TEST} sebagai penyiap \eng{flag} beserta
          perbedaan keduanya.
    \item Lompatan tak bersyarat \asm{JMP} dan jarak lompatan
          \eng{short}, \eng{near}, dan \eng{far}.
    \item Lompatan bersyarat \asm{Jcc} dan pemetaannya terhadap
          bilangan bertanda serta tak bertanda.
    \item Label, aturan penamaan, dan kesalahan umum pada lompatan.
    \item Perulangan: \asm{LOOP}, \asm{LOOPE}/\asm{LOOPZ},
          \asm{LOOPNE}/\asm{LOOPNZ}, dan register \reg{CX}.
    \item Penerjemahan pola bahasa tingkat tinggi ke assembly,
          termasuk tabel lompatan untuk \code{switch}.
    \item Studi kasus: nilai maksimum, pengurutan gelembung, dan
          pengantar pencarian biner.
\end{enumerate}
\end{materibox}

\subfile{section-01}
\subfile{section-02}
\subfile{section-03}

\section{Contoh}
\label{sec:contoh-bab12}

Dua contoh berikut memperlihatkan penerapan mode pengalamatan.
Contoh pertama adalah program utuh yang mengakses array dengan tiga
mode berbeda. Contoh kedua membandingkan ketujuh mode untuk mengakses
elemen yang sama.

\subsection{Contoh 1: Tiga Mode untuk Array yang Sama}

Program pada Kode~\ref{lst:mode-array} menjumlahkan lima elemen array
dengan tiga cara: \textit{direct memory}, \textit{direct index}, dan
\textit{base relative}.

\lstinputlisting[
    language=MASM,
    caption={Program \file{mode-array.asm} membandingkan tiga mode
    pengalamatan.},
    label={lst:mode-array}
]{\CodeDir/mode-array.asm}

Bagian penting program:

\begin{itemize}
    \item Cara pertama memakai \textit{direct memory}:
          \lstasm|MOV AX, data1+0| sampai \lstasm|ADD AX, data1+8|.
          Offset ditulis sebagai konstanta.
    \item Cara kedua memakai \textit{direct index}:
          \lstasm|ADD AX, data1[SI]|. \reg{SI} dinaikkan dua setiap
          elemen karena satu word berukuran dua byte.
    \item Cara ketiga memakai \textit{base relative}:
          \lstasm|ADD AX, [BX+SI]|. \reg{BX} menunjuk awal array dan
          \reg{SI} menjadi indeks.
    \item Ketiga hasil dibandingkan dengan \asm{CMP}. Bila sama,
          program mencetak \code{Ketiga hasil identik.}
\end{itemize}

Keluaran program tampak pada Kode~\ref{lst:keluaran-mode}.

\begin{lstlisting}[
    style=sessiongaya,
    caption={Keluaran program \file{mode-array.asm}.},
    label={lst:keluaran-mode}
]
Perbandingan tiga mode pengalamatan
Ketiga hasil identik.
\end{lstlisting}

\subsection{Contoh 2: Tujuh Mode untuk Satu Elemen}

Misalkan array \code{data1 DW 10,20,30,40,50} dan elemen yang dituju
adalah elemen ke-2 (indeks 2, offset 4, nilai 30). Ketujuh mode dapat
dipakai seperti pada Tabel~\ref{tab:tujuh-mode-contoh}.

\begin{table}[H]
    \centering
    \caption{Tujuh mode pengalamatan untuk mengakses
    \code{data1[2]}.}
    \label{tab:tujuh-mode-contoh}
    \small
    \begin{tabular}{@{}p{2.6cm}p{6.2cm}p{4.2cm}@{}}
        \toprule
        \textbf{Mode} & \textbf{Potongan program} &
        \textbf{Alamat efektif} \\
        \midrule
        \textit{Register} & \lstasm|MOV AX, BX| &
        Tidak ada akses memori \\
        \addlinespace
        \textit{Immediate} & \lstasm|MOV AX, 30| &
        Tidak ada akses memori \\
        \addlinespace
        \textit{Direct} & \lstasm|MOV AX, data1+4| &
        offset \code{data1} $+ 4$ \\
        \addlinespace
        \textit{Indirect} & \lstasm|LEA SI, data1+4| lalu
        \lstasm|MOV AX, [SI]| & offset \code{data1} $+ 4$ \\
        \addlinespace
        \textit{Base relative} & \lstasm|LEA BX, data1| lalu
        \lstasm|MOV AX, [BX+4]| & offset \code{data1} $+ 4$ \\
        \addlinespace
        \textit{Direct index} & \lstasm|LEA SI, data1| lalu
        \lstasm|MOV AX, [SI+4]| & offset \code{data1} $+ 4$ \\
        \addlinespace
        \textit{Based indexed} & \lstasm|LEA BX, data1|,
        \lstasm|MOV SI, 4|, lalu \lstasm|MOV AX, [BX+SI]| &
        offset \code{data1} $+ 4$ \\
        \bottomrule
    \end{tabular}
\end{table}

Kelima mode memori menghasilkan nilai yang sama, yaitu 30, karena
alamat efektifnya sama. Yang berbeda adalah cara membentuk alamat
tersebut. Mode \textit{register} dan \textit{immediate} tidak
mengakses memori sama sekali.

\begin{infobox}
\vspace{1.1em}
Pada mode \textit{direct index}, perpindahan ditulis di dalam
kurung siku bersama register indeks, misalnya \lstasm|[SI+4]|.
Pada mode \textit{indirect}, kurung siku hanya berisi register,
misalnya \lstasm|[SI]|. Perbedaan kecil ini menentukan apakah
alamat efektif ditambah perpindahan atau tidak.
\end{infobox}

\section{Praktikum}
\label{sec:praktikum-bab12}

\begin{praktikumbox}
\vspace{1.1em}
\textbf{Judul:} Membandingkan mode pengalamatan pada akses array.\\[0.3em]
\textbf{Sub-CPMK:} 3.3 \quad\textbullet\quad
\textbf{Alokasi waktu:} 100 menit \quad\textbullet\quad
\textbf{Bentuk:} individu
\end{praktikumbox}

\subsection{Tujuan Praktikum}

Setelah menyelesaikan praktikum ini, mahasiswa mampu menggunakan
tujuh mode pengalamatan, menghitung alamat efektif, dan membuktikan
bahwa beberapa mode dapat menunjuk elemen array yang sama.

\subsection{Alat dan Bahan}

\begin{itemize}
    \item Komputer dengan emu8086 terpasang.
    \item DOSBox dengan TASM, TLINK, dan DEBUG.
    \item Berkas \file{mode-array.asm} dari \file{code/asm/}.
    \item Lembar pengamatan register dan alamat.
\end{itemize}

\subsection{Langkah Kerja}

\begin{enumerate}
    \item \textbf{Merakit dan menjalankan program.}
    \begin{enumerate}
        \item Buka DOSBox, lalu masuk ke direktori kerja.
        \item Rakit dengan \file{tasm mode-array.asm}.
        \item Link dengan \file{tlink mode-array.obj}.
        \item Jalankan program dan catat keluarannya.
    \end{enumerate}

    \item \textbf{Menelusuri tiga mode.}
    \begin{enumerate}
        \item Buka program pada emu8086 dan jalankan
              \eng{single step}.
        \item Catat \reg{AX}, \reg{SI}, dan \reg{BX} pada bagian
              \textit{direct}, \textit{direct index}, dan
              \textit{base relative}.
        \item Hitung alamat efektif setiap akses dengan rumus
              $EA = basis + indeks + displacement$.
        \item Bandingkan nilai \code{jumlah1}, \code{jumlah2}, dan
              \code{jumlah3} pada jendela memori.
    \end{enumerate}

    \item \textbf{Mengubah isi array.}
    \begin{enumerate}
        \item Ubah elemen array menjadi \code{10, 20, 30, 40, 100}.
        \item Rakit ulang dan jalankan program.
        \item Pastikan ketiga cara tetap menghasilkan jumlah yang
              sama.
        \item Jelaskan mengapa perubahan isi tidak mengubah
              hubungan antarmode.
    \end{enumerate}

    \item \textbf{Mencoba mode lain.}
    \begin{enumerate}
        \item Ubah cara kedua dari \lstasm|data1[SI]| menjadi
              \lstasm|[SI+data1]|.
        \item Ubah cara ketiga dari \lstasm|[BX+SI]| menjadi
              \lstasm|[BX+4]| untuk elemen ke-2.
        \item Rakit ulang dan pastikan hasilnya tetap sama.
    \end{enumerate}

    \item \textbf{Memeriksa alamat dengan DEBUG.}
    \begin{enumerate}
        \item Jalankan \file{debug mode-array.exe}.
        \item Catat nilai \reg{DS} dengan perintah \file{r}.
        \item Hitung alamat fisik awal array dengan rumus
              $(\reg{DS} \times \hex{10}) + \text{offset data1}$.
        \item Tampilkan isi array dengan
              \lstasm|D DS:offset L 0A|.
        \item Cocokkan urutan word dengan nilai array.
    \end{enumerate}

    \item \textbf{Menulis potongan tujuh mode.}
    \begin{enumerate}
        \item Tulis potongan program yang mengakses elemen ke-2
              dengan mode \textit{direct}, \textit{indirect},
              \textit{base relative}, \textit{direct index}, dan
              \textit{based indexed}.
        \item Jelaskan mode mana yang tidak mengakses memori.
    \end{enumerate}
\end{enumerate}

\subsection{Hasil yang Dikumpulkan}

\begin{enumerate}
    \item Berkas \file{mode-array.asm} versi modifikasi.
    \item Tangkapan layar keluaran program.
    \item Tabel alamat efektif untuk tiga mode.
    \item Tangkapan layar isi array pada DEBUG.
    \item Penjelasan perbedaan mode \textit{indirect} dan
          \textit{direct index}.
\end{enumerate}

\begin{asesmenbox}
\vspace{1.1em}
\textbf{Asesmen praktikum 10 (bobot 5\% dari komponen praktikum).}

\begin{tabular}{@{}p{6.2cm}cp{4.6cm}@{}}
    \toprule
    \textbf{Indikator} & \textbf{Bobot} & \textbf{Kriteria} \\
    \midrule
    Ketepatan penggunaan tiga mode & 35\% &
    Ketiga cara mengakses elemen dengan benar \\
    \addlinespace
    Ketepatan perhitungan alamat efektif & 30\% &
    Basis, indeks, dan displacement dijumlahkan benar \\
    \addlinespace
    Ketepatan analisis hasil & 20\% &
    Kesamaan hasil dijelaskan melalui alamat efektif \\
    \addlinespace
    Kerapian dokumentasi & 15\% &
    Lembar kerja lengkap dan dapat diulang \\
    \bottomrule
\end{tabular}
\end{asesmenbox}

\section{Latihan}
\label{sec:latihan-bab12}

\begin{enumerate}
    \item \textbf{Konsep.} Jelaskan pengertian mode pengalamatan.
          Sebutkan tujuh mode yang dibahas pada bab ini.

    \item \textbf{Identifikasi.} Tentukan mode pengalamatan pada
          setiap operand memori berikut.
          \begin{enumerate}
              \item \lstasm|MOV AX, BX|
              \item \lstasm|MOV AX, 1234h|
              \item \lstasm|MOV AX, nilai|
              \item \lstasm|MOV AX, [SI]|
              \item \lstasm|MOV AX, [BX+4]|
              \item \lstasm|MOV AX, [DI+8]|
              \item \lstasm|MOV AX, [BX+SI+2]|
          \end{enumerate}

    \item \textbf{Direct dan indirect.} Jelaskan perbedaan mode
          \textit{direct memory} dan \textit{indirect memory}.
          Berikan satu contoh untuk masing-masing.

    \item \textbf{Base relative dan direct index.} Jelaskan
          perbedaan \lstasm|[BX+4]| dan \lstasm|[SI+4]|. Mode apa
          yang dipakai oleh masing-masing bentuk?

    \item \textbf{Based indexed.} Sebutkan empat kombinasi register
          yang legal untuk mode \textit{based indexed}. Mengapa
          \lstasm|[BX+BP]| dan \lstasm|[SI+DI]| tidak legal?

    \item \textbf{Alamat efektif.} Diketahui \reg{DS} = \hex{2000},
          \reg{BX} = \hex{0100}, dan \reg{SI} = \hex{0006}.
          Hitung alamat efektif dan alamat fisik untuk
          \lstasm|MOV AX, [BX+SI+4]|.

    \item \textbf{Segmen default.} Tentukan segmen default untuk
          \lstasm|[BX]|, \lstasm|[BP]|, \lstasm|[SI]|, dan
          \lstasm|[DI]|. Mengapa \lstasm|MOV AX, [BP]| tidak membaca
          variabel pada segmen data?

    \item \textbf{Pemilihan mode.} Tentukan mode yang paling tepat
          untuk:
          \begin{enumerate}
              \item variabel tunggal;
              \item elemen array satu dimensi dengan indeks berubah;
              \item medan record;
              \item elemen array of record.
          \end{enumerate}
\end{enumerate}

\begin{refleksibox}
\vspace{1.1em}
Jawab pertanyaan berikut pada jurnal praktikum Anda.
\begin{enumerate}
    \item Mode mana yang paling sulit Anda bedakan: \textit{indirect},
          \textit{base relative}, atau \textit{direct index}?
          Mengapa?
    \item Setelah menghitung alamat efektif, apakah Anda memahami
          mengapa beberapa penulisan menghasilkan alamat yang sama?
          Jelaskan.
    \item Konsep apa yang perlu Anda ulangi sebelum mempelajari
          percabangan dan perulangan pada bab berikutnya?
\end{enumerate}
\end{refleksibox}

\section{Rangkuman}
\label{sec:rangkuman-bab12}

Pada bab ini telah dibahas:

\begin{itemize}
    \item Pengambilan keputusan berjalan dalam dua langkah: satu
          instruksi menyiapkan \eng{flag}, lalu satu instruksi
          lompatan bersyarat membacanya. Lima \eng{flag} terpenting
          adalah \flag{ZF}, \flag{CF}, \flag{SF}, \flag{OF}, dan
          \flag{PF} (Bab~\ref{chap:register}).

    \item \asm{CMP} menyiapkan \eng{flag} melalui pengurangan yang
          hasilnya dibuang, sedangkan \asm{TEST} melalui \asm{AND}
          yang hasilnya dibuang. \asm{CMP} membandingkan nilai;
          \asm{TEST} memeriksa bit.

    \item \asm{JMP} memindahkan kendali tanpa syarat. Jarak lompatan
          menentukan ukuran instruksi: \eng{short} dua byte
          ($\pm$ 127 byte), \eng{near} tiga byte (satu segmen), dan
          \eng{far} lima byte (antarsegmen).

    \item Keluarga \asm{Jcc} terbagi menjadi kelompok bertanda
          (\asm{JG}, \asm{JGE}, \asm{JL}, \asm{JLE}), kelompok tak
          bertanda (\asm{JA}, \asm{JAE}, \asm{JB}, \asm{JBE}),
          pemeriksa \eng{flag} tunggal (\asm{JC}, \asm{JS},
          \asm{JO}, \asm{JP}, \asm{JCXZ}), dan kelompok kesamaan
          (\asm{JE}/\asm{JNE}) yang berlaku untuk keduanya.

    \item Perulangan dibentuk dengan \asm{LOOP} dan \reg{CX}, varian
          \asm{LOOPE}/\asm{LOOPNE} yang menambahkan syarat \flag{ZF},
          atau pola \asm{DEC} dan \asm{JNZ} yang lebih luwes dan
          lebih cepat pada prosesor modern.

    \item Seluruh struktur kontrol bahasa tingkat tinggi---\code{if},
          \code{while}, \code{for}, dan \code{switch}---diterjemahkan
          dari tiga bahan dasar: \asm{CMP}/\asm{TEST}, \asm{Jcc}, dan
          \asm{JMP}. \code{switch} dengan banyak cabang paling
          efisien memakai tabel lompatan.

    \item Studi kasus pencarian maksimum, pengurutan gelembung, dan
          pencarian biner menegaskan bahwa struktur kontrol kompleks
          tetap dibentuk dari bahan dasar yang sama.
\end{itemize}

\begin{checklistbox}
Tandai setiap butir yang sudah Anda kuasai.

\begin{itemize}
    \item[\checkbox] Saya dapat menjelaskan peran \eng{flag} dan
          menyebutkan instruksi yang menyiapkannya.
    \item[\checkbox] Saya dapat membedakan \asm{CMP} dan \asm{TEST}
          beserta situasi pemakaian masing-masing.
    \item[\checkbox] Saya dapat menjelaskan tiga jenis jarak lompatan
          dan pengaruhnya pada ukuran kode.
    \item[\checkbox] Saya dapat memilih \asm{Jcc} yang benar untuk
          data bertanda dan tak bertanda tanpa melihat tabel.
    \item[\checkbox] Saya dapat menulis perulangan dengan \asm{LOOP}
          maupun dengan \asm{DEC} dan \asm{JNZ}.
    \item[\checkbox] Saya dapat menerjemahkan pseudocode
          \code{if--else} dan \code{while} ke assembly.
    \item[\checkbox] Saya dapat menjelaskan cara kerja tabel lompatan
          pada pola \code{switch}.
    \item[\checkbox] Saya dapat menelusuri alur percabangan dengan
          \eng{Single Step} dan menjelaskan keputusan setiap
          lompatan.
\end{itemize}

\vspace{0.4em}
Bila masih ada butir yang belum bertanda, ulangi bagian yang
bersesuaian sebelum melanjutkan ke bab berikutnya.
\end{checklistbox}

\begin{infobox}
\textbf{Rujukan bab.} Materi pada bab ini diperkuat oleh
\cite{bartlett2004}, \cite{hyde2010}, \cite{evans2020}, dan
\cite{irvine2011}. Untuk penelusuran lebih dalam mengenai perulangan
dan pemilihan instruksi lompatan pada program nyata, lihat
\cite{irvine2011} dan \cite{hyde2010}. Penerapan tabel lompatan pada
program besar dibahas pada \cite{evans2020}.
\end{infobox}

\section{Evaluasi}
\label{sec:evaluasi-bab12}

\subsection{Soal Pilihan Ganda}

\begin{enumerate}
    \item \lstasm|MOV AX, [BX]| memakai mode pengalamatan \dots
          \begin{enumerate}
              \item \textit{register}
              \item \textit{direct memory}
              \item \textit{indirect memory}
              \item \textit{immediate}
          \end{enumerate}

    \item \lstasm|MOV AX, [BX+SI+4]| memakai mode \dots
          \begin{enumerate}
              \item \textit{base relative}
              \item \textit{direct index}
              \item \textit{based indexed}
              \item \textit{indirect}
          \end{enumerate}

    \item Segmen default untuk \lstasm|MOV AX, [BP]| adalah \dots
          \begin{enumerate}
              \item \reg{CS} \item \reg{DS}
              \item \reg{SS} \item \reg{ES}
          \end{enumerate}

    \item Diketahui \reg{DS} = \hex{1000}, \reg{BX} = \hex{0200},
          \reg{SI} = \hex{0004}, dan \lstasm|MOV AX, [BX+SI+2]|.
          Alamat efektifnya adalah \dots
          \begin{enumerate}
              \item \hex{0204} \item \hex{0206}
              \item \hex{10206} \item \hex{10204}
          \end{enumerate}

    \item Kombinasi register yang legal untuk mode \textit{based
          indexed} adalah \dots
          \begin{enumerate}
              \item \code{[BX+BP]}
              \item \code{[SI+DI]}
              \item \code{[BX+DI]}
              \item \code{[AX+SI]}
          \end{enumerate}
\end{enumerate}

\subsection{Soal Uraian}

\begin{enumerate}
    \item Sebutkan tujuh mode pengalamatan 8086 dan berikan satu
          contoh untuk masing-masing.

    \item Jelaskan perbedaan mode \textit{direct memory},
          \textit{indirect memory}, dan \textit{base relative}.
          Berikan contoh penulisan untuk setiap mode.

    \item Sebutkan empat kombinasi register yang legal pada mode
          \textit{based indexed}. Hitung alamat efektif dan alamat
          fisik untuk \lstasm|MOV AX, [BX+SI+4]| bila
          \reg{DS} = \hex{2000}, \reg{BX} = \hex{0100}, dan
          \reg{SI} = \hex{0006}.

    \item Jelaskan segmen default untuk \reg{BX}, \reg{BP},
          \reg{SI}, dan \reg{DI}. Tentukan pula mode yang tepat untuk
          mengakses variabel tunggal, elemen array, medan record, dan
          elemen array of record.
\end{enumerate}

\subsection{Kisi-kisi Penilaian}

\begin{table}[H]
    \centering
    \caption{Kisi-kisi evaluasi Bab~\ref{chap:mode-pengalamatan}.}
    \label{tab:kisi-bab12}
    \small
    \begin{tabular}{@{}p{1.5cm}p{5.0cm}p{3.0cm}c@{}}
        \toprule
        \textbf{Sub-CPMK} & \textbf{Indikator yang diukur} &
        \textbf{Bentuk soal} & \textbf{Bobot} \\
        \midrule
        3.3 & Mengidentifikasi tujuh mode pengalamatan &
        Pilihan ganda 1, 2; Uraian 1, 2 & 35\% \\
        \addlinespace
        3.3 & Menerapkan kombinasi register dan menghitung alamat &
        Pilihan ganda 4; Uraian 3 & 30\% \\
        \addlinespace
        3.3 & Menentukan segmen default &
        Pilihan ganda 3; Uraian 4 & 20\% \\
        \addlinespace
        3.3 & Memilih mode untuk array dan record &
        Pilihan ganda 5; Uraian 4 & 15\% \\
        \bottomrule
    \end{tabular}
\end{table}

\begin{checklistbox}
\vspace{1.1em}
\textbf{Checklist Pencapaian Kompetensi
Bab~\ref{chap:mode-pengalamatan}}

Seorang mahasiswa dinyatakan kompeten pada Sub-CPMK 3.3 apabila:
\begin{itemize}
    \item[\checked] dapat menyebutkan tujuh mode pengalamatan;
    \item[\checked] dapat mengenali mode dari penulisan operand;
    \item[\checked] dapat membedakan \textit{direct} dan
          \textit{indirect};
    \item[\checked] dapat menggunakan \textit{base relative} dan
          \textit{direct index};
    \item[\checked] dapat menggunakan \textit{based indexed};
    \item[\checked] dapat menghitung alamat efektif dan alamat fisik;
    \item[\checked] dapat menentukan segmen default.
\end{itemize}
\end{checklistbox}


\end{document}
